% Pattern: number/wrapper. A wrapper sets \theinner to its argument and opens
% the inner environment, so each use prints the number it gives. Seen in the
% gold pilot (set X).
% want customthm-1 statement/theorem A
% want customthm-2 statement/theorem B
% want theorem-1 statement/theorem 1.1
\documentclass{article}
\usepackage{amsthm}
\newtheorem{theorem}{Theorem}[section]
\newtheorem{innercustomthm}{Theorem}
\newenvironment{customthm}[1]
  {\renewcommand\theinnercustomthm{#1}\innercustomthm}
  {\endinnercustomthm}
\begin{document}
\begin{customthm}{A} Every compact space is bounded. \end{customthm}
\begin{customthm}{B} Every finite space is bounded. \end{customthm}
\section{One}
\begin{theorem} Every finite space is compact. \end{theorem}
\end{document}
